\scanpage{68}
Le foncteur
\[
G\longmapsto X
\]
s'explicite comme
\[
G\longmapsto\underline{\operatorname{Bor}}(G)\tag*{(2\unsure{5|4})}
\]
associant à $G$ le \deleted{\unsure{ens}} schém. de ses sous-groupes de Borel $B$ (i.e. \deleted{\unsure{tels}} des ss-gpes tels que \deleted{$(B,G)$} \inserted{$(G,B)$} \unsure{sont|soit} loc. isom. \deleted{\illegible{2}} (\inserted{\unsure{Zar}}, ou $\mathrm{fpqc}$) à \deleted{$\unsure{B|(G_0}$} $(G_{0S},B_{0S})$. On obtient en fait, si $G$ et $X$ se correspondent -- donc $G\simeq\underline{\operatorname{Aut}}(X)$ -- un isom. can.
\[
X\xrightarrow{\sim}\underline{\operatorname{Bor}}(G)\tag*{(2\unsure{6|5})}
\]
en associant à tte section \inserted{\unsure{locale|loc.}} de $X$ son stabilisateur dans $G$.

\uline{V} Les foncteurs vers les formes de $\mathrm{SL}(2)_S$ s'\deleted{\unsure{fac|f}}obtiennent soit en factorisant par $\operatorname{Tor}(\mathrm{GP}(1)_S)$
\[
\mathcal X\longmapsto R=\underline{\operatorname{Isom}}(\mathcal X_{0S},\mathcal X)\longmapsto R\times_{\mathrm{GP}(1)_S}\unsure{\widetilde G_0|\widetilde G_{0S}},\tag*{(27)}
\]
soit plus \unsure{naturellt|naturellement} en
\[
\mathcal X\longmapsto G=\underline{\operatorname{Aut}}_S(\mathcal X)\longrightarrow\widetilde G=\text{Rev. \unsure{simplemt|simpl. conn.} de $G$}\tag*{(28)}
\]
% archive sheet number: 67
